\begin{figure}[H]
\centering
\begin{tikzpicture}[ancho/.style={minimum width=3.2cm}]
    \node[servicio, ancho] (blogservice) at (0,0) {Servicio de\\publicaciones};
    \node[servicio, ancho] (mediaservice) at (0,-2.1) {Servicio de\\imágenes};
    \node[servicio, ancho] (snapshotservice) at (0,-4.2) {Servicio de\\respaldos};

    \node[documento, minimum width=3.4cm, minimum height=5.3cm] (contentstore) at (6,-2.1) {Content Store\\JSON, Markdown\\e imágenes};
    \node[datos, minimum width=3cm] (sqlite) at (6,-6.7) {SQLite\\operativo};

    \node[servicio, ancho] (publishservice) at (11.4,-2.1) {Servicio de\\publicación};
    \node[interfaz, ancho] (site) at (11.4,-6.7) {Sitio público\\publicado};

    \draw[flecha] (blogservice) -- node[etq, above=2pt] {crea y edita} (blogservice-|contentstore.west);
    \draw[flecha] (mediaservice) -- node[etq, above=2pt] {valida y\\guarda} (mediaservice-|contentstore.west);
    \draw[flecha] (snapshotservice) -- node[etq, above=2pt] {copia el\\contenido} (snapshotservice-|contentstore.west);
    \draw[flecha] (snapshotservice.south) |- node[etq, pos=0.75, above=2pt] {registra metadatos} (sqlite.west);
    \draw[flecha] (contentstore.east|-publishservice) -- node[etq, above=2pt] {datos\\publicables} (publishservice);
    \draw[flecha] (publishservice) -- node[etq, right=3pt] {construye\\el sitio} (site);
    \draw[flecha] (publishservice.south west) -- node[etq, sloped, above=2pt] {registra el estado} (sqlite.east);
\end{tikzpicture}
\caption[Servicios y almacenes del CMS]{Relación entre servicios del CMS y almacenes de contenido, recursos y operación}
\label{fig:cms_servicios_almacenes}
\end{figure}
